\documentclass[11pt,a4paper]{article}

\def\courseTitle{Industrial Security (Master)}
\def\sectionNumber{01}
\def\sectionTitle{Firmware Reverse Engineering}
\def\sheetKind{Exercise Sheet}

% =====================================================================
% Shared preamble for exercise sheets and solution sheets.
%
% Define BEFORE \input{}-ing this file:
%   \def\courseTitle{...}
%   \def\sectionNumber{...}
%   \def\sectionTitle{...}
%   \def\sheetKind{Exercise Sheet}   or  Solution Sheet
%
% Optional:
%   \def\courseAuthor{...}
%
% The caller must issue \documentclass{...} (typically article) BEFORE
% loading this preamble.
% =====================================================================

\usepackage[utf8]{inputenc}
\usepackage[T1]{fontenc}
\usepackage[english]{babel}
\usepackage[a4paper, margin=2.2cm]{geometry}
\usepackage{graphicx}
\usepackage{listings}
\usepackage{xcolor}
\usepackage{tikz}
\usepackage{booktabs}
\usepackage{array}
\usepackage{enumitem}
\usepackage{titlesec}
\usepackage{fancyhdr}
\usepackage{hyperref}
\usepackage{amsmath}
\usepackage{amssymb}
\usepackage{tcolorbox}
\tcbuselibrary{skins, breakable}
\usepackage{needspace}

% Discourage solo lines split across pages.
\widowpenalty=10000
\clubpenalty=10000
\raggedbottom

% =====================================================================
% Shared TikZ styles for the Industrial Security lecture series.
% Loaded by every slide deck and exercise sheet that uses TikZ.
% =====================================================================

\usetikzlibrary{
  arrows.meta,
  shapes,
  shapes.geometric,
  shapes.symbols,
  shapes.multipart,
  positioning,
  fit,
  calc,
  backgrounds,
  decorations.pathmorphing,
  decorations.pathreplacing,
  decorations.text,
  patterns,
  matrix,
  shadows
}

% --- colour palette --------------------------------------------------
\definecolor{isOT}{HTML}{1F4E79}     % OT (deep blue)
\definecolor{isIT}{HTML}{6F2DA8}     % IT (purple)
\definecolor{isSafety}{HTML}{C00000} % safety red
\definecolor{isNet}{HTML}{2E7D32}    % network green
\definecolor{isAttack}{HTML}{B23A48} % attack red
\definecolor{isAccent}{HTML}{E08E0B} % amber
\definecolor{isMute}{HTML}{6B6B6B}   % grey
\definecolor{isLight}{HTML}{EDF2F8}  % light background
\definecolor{isPLC}{HTML}{2C5282}
\definecolor{isHMI}{HTML}{2F855A}
\definecolor{isSCADA}{HTML}{6B46C1}

% --- generic node styles --------------------------------------------
\tikzset{
  isnode/.style={
    draw, rounded corners=2pt, minimum height=8mm, minimum width=18mm,
    font=\small, align=center, thick
  },
  isbox/.style={isnode, fill=isLight},
  plc/.style={isnode, fill=isPLC!15, draw=isPLC, thick},
  hmi/.style={isnode, fill=isHMI!15, draw=isHMI, thick},
  scada/.style={isnode, fill=isSCADA!15, draw=isSCADA, thick},
  sensor/.style={isnode, fill=isNet!10, draw=isNet, minimum height=6mm, minimum width=14mm, font=\scriptsize},
  actuator/.style={isnode, fill=isAccent!15, draw=isAccent, minimum height=6mm, minimum width=14mm, font=\scriptsize},
  firewall/.style={isnode, fill=isAttack!15, draw=isAttack, thick},
  server/.style={isnode, fill=isMute!15, draw=isMute!80, thick},
  switch/.style={isnode, fill=isLight, draw=black!60, minimum height=6mm, minimum width=14mm, font=\scriptsize},
  workstation/.style={isnode, fill=white, draw=isOT, thick},
  attacker/.style={isnode, fill=isAttack!20, draw=isAttack, thick, font=\small\bfseries},
  inetnode/.style={draw, ellipse, minimum width=24mm, minimum height=10mm, font=\scriptsize, fill=isLight, thick},
  process/.style={isnode, fill=isOT!15, draw=isOT, thick, minimum width=24mm},
  zone/.style={
    draw, rounded corners=4pt, thick, dashed, inner sep=4mm
  },
  conduit/.style={-{Stealth[length=2.5mm]}, thick, isNet!80!black},
  attack/.style={-{Stealth[length=2.5mm]}, thick, isAttack, dashed},
  flow/.style={-{Stealth[length=2.5mm]}, thick, isOT!70!black},
  netline/.style={thick, black!60},
  axisline/.style={-{Stealth[length=2mm]}, thick, black!70},
  tag/.style={font=\scriptsize\itshape, fill=white, inner sep=1pt},
  level/.style={
    draw, fill=isLight, minimum width=0.85\linewidth, minimum height=8mm,
    font=\small, anchor=west
  },
  brace/.style={decorate, decoration={brace, amplitude=4pt, raise=2pt}},
  legendbox/.style={draw, rounded corners=2pt, fill=isLight, inner sep=2mm, font=\scriptsize, align=left}
}

% --- handy macros ----------------------------------------------------
% \plcicon, \hmiicon etc as simple labelled boxes; useful inline.
\newcommand{\plcicon}[1][PLC]{\tikz[baseline=-0.6ex]{\node[plc, minimum width=10mm, minimum height=5mm, font=\scriptsize, inner sep=1pt]{#1};}}
\newcommand{\hmiicon}[1][HMI]{\tikz[baseline=-0.6ex]{\node[hmi, minimum width=10mm, minimum height=5mm, font=\scriptsize, inner sep=1pt]{#1};}}
\newcommand{\fwicon}[1][FW]{\tikz[baseline=-0.6ex]{\node[firewall, minimum width=10mm, minimum height=5mm, font=\scriptsize, inner sep=1pt]{#1};}}


\providecommand{\courseAuthor}{Industrial Security Lecture Series}
\providecommand{\courseInstitute}{Department of Computer Science}
\providecommand{\companyLogo}{logo}

% --- Corporate branding overrides ------------------------------------
% Picks up logo / author / brand colours from the repo-root branding.tex
% when present; falls back to the defaults above otherwise.
\IfFileExists{../../../branding.tex}{% =====================================================================
% Corporate / institutional branding -- single file to customise the
% look of the entire course (slides, notes, exercises, labs).
%
% HOW TO REBRAND IN UNDER A MINUTE:
%   1. Replace `branding/logo.pdf` with your own logo
%      (PDF preferred; PNG and JPG also work -- name it `logo.<ext>`).
%   2. (Optional) change the two colours below to your brand palette.
%   3. (Optional) override the author/institute lines below.
%   4. Re-run `make` at the repo root. That's it.
%
% Everything in this file has sensible defaults; you can delete a line
% to fall back to the built-in defaults, or comment it out.
%
% This file is loaded automatically by the slides, exercise, and lab
% preambles -- you never need to \input it manually.
% =====================================================================

% --- Logo --------------------------------------------------------------
% Path is resolved from each leaf-build directory; the default points at
% the shared `branding/` folder at the repo root. If you put your logo
% somewhere else, set the full or relative path here (without extension):
\renewcommand{\companyLogo}{../../../branding/logo}

% --- Author / institute (appears on the title page footer) ------------
\renewcommand{\courseAuthor}{}
\renewcommand{\courseInstitute}{}

% --- Brand colours ----------------------------------------------------
% slideAccent      = primary brand colour (title bands, accent rule)
% slideSecondary   = secondary brand colour (section badge, highlight)
% Both use HTML hex codes. Keep enough contrast against white text.
\definecolor{slideAccent}{HTML}{00205B}      % deep navy
\definecolor{slideSecondary}{HTML}{00A0DC}   % light blue accent
}{}

% --- listings -------------------------------------------------------
\definecolor{codebg}{rgb}{0.96,0.96,0.96}
\lstset{
  backgroundcolor=\color{codebg},
  basicstyle=\ttfamily\footnotesize,
  breaklines=true,
  frame=single,
  numbers=left,
  numberstyle=\tiny\color{black!50},
  showstringspaces=false,
  tabsize=2
}

% --- header / footer ------------------------------------------------
\pagestyle{fancy}
\fancyhf{}
\renewcommand{\headrulewidth}{0.4pt}
\renewcommand{\footrulewidth}{0pt}
\lhead{\small \courseTitle}
\rhead{\small Section \sectionNumber{} -- \sheetKind}
\cfoot{\thepage}

% --- task / solution boxes -----------------------------------------
% `before={...}` guards every task/solution box with \needspace so the
% box doesn't start with only one or two lines of breathing room at the
% bottom of a page. If less than ~9 lines remain, LaTeX bumps the box
% to the next page; otherwise it starts here and breaks normally if it
% runs long.
% Boxes are `breakable` so very long solutions don't blow the page,
% but `before={\needspace{14\baselineskip}}` pushes them to a fresh
% page whenever fewer than ~14 lines remain. That covers the vast
% majority of single-question splits in practice.
\newtcolorbox{taskbox}[1]{
  enhanced, breakable,
  colback=isLight, colframe=isOT,
  fonttitle=\bfseries\color{white},
  coltitle=white,
  colbacktitle=isOT,
  title={#1},
  arc=2pt, boxrule=0.6pt,
  left=2mm, right=2mm, top=2mm, bottom=2mm,
  before={\par\vspace{2mm}\needspace{14\baselineskip}\par},
  after={\par\vspace{2mm}}
}

\newtcolorbox{solutionbox}{
  enhanced, breakable,
  colback=isNet!5, colframe=isNet,
  fonttitle=\bfseries\color{white},
  coltitle=white,
  colbacktitle=isNet,
  title={Solution},
  arc=2pt, boxrule=0.6pt,
  left=2mm, right=2mm, top=2mm, bottom=2mm,
  before={\par\vspace{2mm}\needspace{14\baselineskip}\par},
  after={\par\vspace{2mm}}
}

\newtcolorbox{hintbox}{
  enhanced, breakable,
  colback=isAccent!10, colframe=isAccent,
  fonttitle=\bfseries\color{white}, coltitle=white, colbacktitle=isAccent,
  title={Hint},
  arc=2pt, boxrule=0.6pt
}

\newtcolorbox{warnbox}{
  enhanced, breakable,
  colback=isAttack!8, colframe=isAttack,
  fonttitle=\bfseries\color{white}, coltitle=white, colbacktitle=isAttack,
  title={Caution},
  arc=2pt, boxrule=0.6pt
}

% --- section formatting --------------------------------------------
\titleformat{\section}{\Large\bfseries\color{isOT}}{\thesection}{0.5em}{}
\titleformat{\subsection}{\large\bfseries\color{isOT!80!black}}{\thesubsection}{0.5em}{}

% --- title block macro ----------------------------------------------
\newcommand{\sheetheader}{%
  \begin{center}
    {\Large\bfseries \courseTitle}\\[0.3em]
    {\large \sheetKind{} -- Section \sectionNumber}\\[0.2em]
    {\large \sectionTitle}\\[0.2em]
    \rule{\linewidth}{0.4pt}
  \end{center}
  \vspace{0.5em}
}


\begin{document}
\sheetheader

\noindent
\textbf{Goal.} These exercises are pen-and-analysis counterparts to the hands-on lab. They
assume you can already drive \texttt{binwalk}, a SquashFS/JFFS2 extractor, \texttt{strings},
Ghidra, and \texttt{qemu-user}; here you \emph{reason} about the artefacts those tools produce
and about the disclosure and supply-chain process around a finding. Answers should be concise
(3--10 sentences each unless stated otherwise), grounded in the cited standards, advisories,
and CWE catalogue, and framed for \emph{authorised} analysis only. Several questions give you a
realistic artefact (an entropy curve, a decompiler listing, a version diff) and ask you to draw
the conclusion a PSIRT or red-team analyst would draw.

\begin{taskbox}{Exercise 1.1 -- Lawful Sourcing and Chain of Custody}
You receive a firmware image three ways: (a)~a vendor support-portal download you are entitled
to under a maintenance contract, (b)~a SPI-flash dump you read from a device you own, and
(c)~an image a colleague forwarded ``from a customer site''. For each:
\begin{enumerate}\setlength{\itemsep}{0pt}
  \item State whether you may lawfully analyse it, and on what legal basis (cite the relevant
        EU and German/US provisions from the slides).
  \item List the minimum chain-of-custody metadata you must record before any analysis, and
        explain in one sentence why a missing element makes a later disclosure worthless.
  \item Explain the difference between a vendor-published \emph{hash} and a detached
        \emph{signature}, and which one actually defends against a tampered download mirror.
\end{enumerate}
\end{taskbox}

\begin{taskbox}{Exercise 1.2 -- Architecture and Endianness Triage}
\texttt{binwalk -A} on an unknown image reports a dense cluster of \texttt{MIPS instructions}
hits but also a thinner scatter of \texttt{ARM} and \texttt{PPC} hits; \texttt{file} on the
carved \texttt{vmlinux} says nothing useful; the \texttt{uImage} header is absent (the carve
started mid-image).
\begin{enumerate}\setlength{\itemsep}{0pt}
  \item Explain why an opcode-frequency scanner produces \emph{some} hits for the wrong ISA,
        and how you decide which architecture is real. Name two independent cross-checks.
  \item You suspect MIPS but disassembly is garbage. How do endianness (MIPS BE vs.\ LE) and
        a wrong load/base address each corrupt the listing, and how do you tell the two
        failure modes apart from the symptoms?
  \item Given the architecture, what is the \emph{single} most consequential parameter you must
        set correctly before Ghidra cross-references become trustworthy, and why?
\end{enumerate}
\end{taskbox}

\begin{taskbox}{Exercise 1.3 -- Reading an Entropy Curve}
\texttt{binwalk -E} produces the entropy plot sketched below (Shannon entropy on $y\in[0,1]$
against file offset). Region~A sits flat at $\approx 0.45$; region~B ramps to $\approx 0.92$
and is flagged by \texttt{binwalk} as \texttt{LZMA compressed data}; region~C sits flat at
$\approx 0.99$ with \emph{no} compression signature and a sharp square edge at both ends.
\begin{center}
\begin{tikzpicture}[x=2.0mm, y=18mm]
  \draw[->] (0,0) -- (62,0) node[right, font=\tiny]{offset};
  \draw[->] (0,0) -- (0,1.15) node[above, font=\tiny]{entropy};
  \draw[thick] (0,0.45) -- (18,0.45);
  \draw[thick] (18,0.45) -- (24,0.92) -- (40,0.92);
  \draw[thick] (40,0.92) -- (40,0.99) -- (58,0.99) -- (58,0.30);
  \node[font=\tiny] at (9,0.55) {A};
  \node[font=\tiny] at (32,1.02) {B};
  \node[font=\tiny] at (49,1.08) {C};
\end{tikzpicture}
\end{center}
\begin{enumerate}\setlength{\itemsep}{0pt}
  \item Classify each region (plaintext/code, compressed, encrypted or packed) and justify
        each call from its entropy level \emph{and} its shape.
  \item Why can entropy alone never distinguish ``compressed'' from ``encrypted'', and what
        single additional observation breaks the tie for region~C?
  \item If region~C is genuinely encrypted, where in the image do you pivot to recover the
        plaintext, and why is fighting the blob directly the wrong move?
\end{enumerate}
\end{taskbox}

\begin{taskbox}{Exercise 1.4 -- Hard-coded Secret, No Patch Available}
Triage of an end-of-life HMI surfaces three secrets: (i)~\texttt{root:\$1\$ab\$...} (an MD5-crypt
hash) in \texttt{/etc/passwd}; (ii)~a PEM private key in \texttt{/etc/ssl/private/server.key}
that is byte-identical across every unit of the model; (iii)~a literal
\texttt{SUPER\_PASS=Zte521} compiled into the web daemon. The vendor is out of business; no
firmware update will ever ship.
\begin{enumerate}\setlength{\itemsep}{0pt}
  \item Assign the correct CWE to each of the three findings.
  \item For each, propose a \emph{compensating control} the asset owner can deploy
        \emph{without} a firmware patch (network, monitoring, or process). Be specific -- name
        the control, not ``segment the network''.
  \item Rank the three by residual risk after your controls are in place, and justify the
        ranking in one sentence.
\end{enumerate}
\end{taskbox}

\begin{taskbox}{Exercise 1.5 -- Reading a Ghidra Decompilation}
The decompiler produces the following for the OTA update handler of a routable ICS gateway
(ARM LE, user-space daemon). Read it as the artefact, not as code you wrote:
\begin{lstlisting}[language=C]
int apply_update(char *url, char *sig_hdr) {
  char path[64];
  download(url, "/tmp/fw.bin");
  sprintf(path, "/tmp/%s", basename_of(url));   /* (1) */
  if (crc32_file("/tmp/fw.bin") == header_crc(sig_hdr))  /* (2) */
    return flash_write("/tmp/fw.bin");          /* (3) */
  return -1;
}
\end{lstlisting}
\begin{enumerate}\setlength{\itemsep}{0pt}
  \item Identify \emph{two} distinct weaknesses, each with its CWE identifier, and point to the
        marked line for each.
  \item Which weakness is more severe for an Internet-exposed gateway, and why? Argue from the
        attacker's required position, not just impact.
  \item Sketch the patch for the more severe weakness at the level of ``name the function and
        the check'' -- vague ``add validation'' earns nothing.
\end{enumerate}
\end{taskbox}

\begin{taskbox}{Exercise 1.6 -- Static vs.\ Emulation: Choosing the Tool}
You have confirmed a suspicious compare in Ghidra but want dynamic confirmation. For each of the
following targets, decide whether \texttt{qemu-user}, full-system \texttt{qemu-system}, or
hardware-in-the-loop (e.g.\ Avatar$^2$) is the right next step, and state the \emph{one} property
of the target that forces that choice:
\begin{enumerate}\setlength{\itemsep}{0pt}
  \item A statically-linked CGI binary whose only inputs are environment variables and stdin.
  \item A web daemon that reads model/serial from an NVRAM \texttt{ioctl} on \texttt{/dev/mtd3}
        at start-up and refuses to run otherwise.
  \item A bare-metal VxWorks RTOS image with a memory-mapped peripheral the auth routine polls.
\end{enumerate}
For case~(2), name two concrete techniques to get the daemon past the NVRAM read \emph{without}
the device.
\end{taskbox}

\begin{taskbox}{Exercise 1.7 -- Version Diff and the Silent Fix}
A \texttt{diff -rq rootfs.old rootfs.new} between v1.4 and v1.5 of one model returns exactly two
differing files: \texttt{/usr/sbin/httpd} and \texttt{/etc/version}. A \texttt{radiff2 -C} on
\texttt{httpd} shows one changed function, \texttt{check\_login}, where a
\texttt{strcmp(pass,"Zte521")} branch present in v1.4 is gone in v1.5. The vendor's changelog
says only ``stability improvements''.
\begin{enumerate}\setlength{\itemsep}{0pt}
  \item Explain why this two-file, one-function delta is strong evidence of a deliberately
        \emph{silent} security fix rather than a routine change.
  \item Everyone still on v1.4 now has a known-exploitable backdoor. Walk through the ethics:
        is the v1.4 backdoor a reportable finding, and to whom -- the now-fixed vendor, a CNA,
        or the public? What changes if the vendor refuses a CVE?
  \item Why is the silently-fixed branch an \emph{n-day} rather than a 0-day, and why does that
        distinction matter for an asset owner's patch-prioritisation decision?
\end{enumerate}
\end{taskbox}

\begin{taskbox}{Exercise 1.8 -- Draft a Coordinated-Disclosure Email}
Write the actual email you would send to a vendor PSIRT for the OTA weakness of Exercise~1.5
(treat the product as ``Acme IG-200 industrial gateway, firmware v3.1,
SHA-256 \texttt{3f9a...c21e}''). Your email must contain, and only contain:
\begin{itemize}\setlength{\itemsep}{0pt}
  \item Affected product, firmware version, and image hash.
  \item A factual technical summary and \emph{minimal} reproduction -- no working exploit code.
  \item A proposed disclosure timeline (vendor / national CERT / public) with concrete day
        counts appropriate to ICS rollouts.
  \item Your CWE classification and a suggested remediation.
  \item Contact details and a PGP fingerprint placeholder; an explicit authorised-use statement.
\end{itemize}
Keep it under one page. The marker will deduct for threats, deadlines framed as ultimatums,
exploit code, or any claim you cannot substantiate from the artefact.
\end{taskbox}

\begin{taskbox}{Exercise 1.9 -- Supply Chain and SBOM (Short Essay)}
The Boa web server (abandoned upstream since 2005) was found embedded in countless ICS HMIs in
2022; CODESYS runtime CVEs (\emph{Code16}, 2023) cascaded across hundreds of products from many
vendors. In about 350--500 words, argue how a \emph{Software Bill of Materials} (SBOM) -- e.g.\
SPDX or CycloneDX -- changes the asset owner's position when such a shared-component flaw drops.
Your essay must:
\begin{enumerate}\setlength{\itemsep}{0pt}
  \item Explain concretely what an SBOM lets an asset owner do on the day a CVE lands in a
        third-party component that they could \emph{not} do without one.
  \item Name two hard limits of SBOMs (think: transitive dependencies, firmware that ships no
        SBOM, version-pinning vs.\ actual binary content) and why firmware RE remains necessary
        even where SBOMs exist.
  \item Tie the argument to one regulatory driver (US EO~14028 / NTIA minimum elements, or the
        EU Cyber Resilience Act) and state one obligation it places on the \emph{vendor}.
\end{enumerate}
\end{taskbox}

\begin{warnbox}
Authorised use only. Every artefact above is hypothetical or drawn from public advisories.
Analyse only firmware you may lawfully possess, on hardware you own; never run a derived finding
(credential, backdoor, or exploit) against a device you do not own or are not written-authorised
to test. Keep all proof-of-concept private until coordinated disclosure with the vendor PSIRT
and the relevant national CERT is complete.
\end{warnbox}

\end{document}
